\documentclass[12pt]{article}
\usepackage[utf8]{inputenc}
\usepackage[T1]{fontenc}
\usepackage{hyperref}
\title{Wano Language Resources for K-12 Curriculum Development}
\author{Compiled by Claude Sonnet 5.5 (claude-sonnet-5-5) for Andr\'as Kornai}
\date{}
\begin{document}
\maketitle

\section{Introduction}
Wano (ISO 639-3: wno; Glottocode wano1243) is an endangered Trans-New Guinea language of Papua, Indonesia. It is spoken in the central highlands, in the Puncak and Puncak Jaya regencies. It is closely related to Dani, Walak and Nggem. Speaker estimates vary widely. Wikipedia gives about 1,000 (2011), Joshua Project gives 1,200, and older SIL-based figures are near 7,000. The Endangered Languages Project says it is under pressure from Western Dani, is losing its trade-language role, and has some literacy.

\section{Primary Educational Resources}
No primers, textbooks or school programmes in Wano were found in the searches. Burung's work is the main scholarly source, and no teaching materials came from it.

\section{Linguistic and Reference Materials}
\begin{itemize}
\item Wikipedia entry on the Wano language: \url{https://en.wikipedia.org/wiki/Wano_language}
\item Burung, abstract on Wano (Max Planck Institute for Evolutionary Anthropology): \url{https://lingweb.eva.mpg.de/archive/wlp/3/abstracts/Burung.pdf}
\item LINGUIST List issue 29.2656: \url{https://linguistlist.org/issues/29/2656}
\end{itemize}

\section{Religious and Community Resources}
Joshua Project has a profile of the Wano of Indonesia: \url{https://joshuaproject.net/people_groups/15857}. No scripture translations were found in the searches.

\section{Digital and Audio Resources}
\begin{itemize}
\item Endangered Languages Project entry for Wano: \url{https://endangeredlanguages.com/lang/2800}
\item DoBeS archive page (MPI Nijmegen) that appeared in the search results: \url{https://dobes.mpi.nl/?p=1034}
\end{itemize}
I did not confirm what recordings the DoBeS page holds, or whether they relate to Wano.

\section{Curriculum Development Context}
Wano is used as a first language by some people in every generation. It is not known to be taught in schools. Any curriculum would have to start from the descriptive work above. It would also need to take account of the dominance of Western Dani in the region.
\end{document}